\chapter{Publication Output}
\label{ch:output}

\section{Still images}

\menu{File}\menusep\menu{Import / Export}\menusep\menu{Export Image}
(\kbd{Ctrl}+\kbd{E}) renders the current view off-screen at any resolution,
independent of the window size, with 8$\times$ multisample anti-aliasing.

\begin{description}
  \item[\ui{Resolution preset}] 720p, 1080p, 4K, or custom width and
    height up to 8192\,px.
  \item[\ui{Background}] \ui{Transparent (PNG only)}, \ui{Solid white}, or
    a custom colour.
  \item[Format] PNG or JPEG, chosen by the extension.
\end{description}

\begin{caltip}
Transparent PNG is the right choice for figures that will sit on a coloured
background in a paper or slide. Combine it with an orthographic, axis
aligned view (\kbd{O}, then an alignment button) for a clean, reproducible
result.
\end{caltip}

\section{Animations}

\menu{File}\menusep\menu{Import / Export}\menusep\menu{Export Animation
(GIF/MP4)} writes either a turntable rotation of a static structure or a
playback of a trajectory.

\begin{description}
  \item[Source] \ui{Turntable rotation (360\textdegree)} or the current
    trajectory.
  \item[\ui{Resolution preset}] Up to 4K.
  \item[\ui{Frames per second}] Playback rate.
  \item[\ui{Background}] Including \ui{Transparent (GIF only)}.
  \item[Format] Animated GIF or H.264 MP4.
\end{description}

\begin{calwarn}
GIF export needs \texttt{pillow}; MP4 needs \texttt{imageio} and
\texttt{imageio-ffmpeg}. MP4 has no alpha channel --- asking for a
transparent MP4 falls back to a solid background.
\end{calwarn}

\section{Ray-traced rendering}

\menu{File}\menusep\menu{Import / Export}\menusep\menu{Ray-Traced Render}
exports the active scene --- camera, lights, representation, colours,
multiple bonds, unit cell --- to an external ray tracer for
publication-quality output with true shadows and reflections.

\begin{description}
  \item[\ui{Engine}] \ui{POV-Ray} or \ui{Tachyon}.
  \item[\ui{Width (px)}, \ui{Height (px)}] Default $1920 \times 1440$.
  \item[\ui{Background}] \ui{Solid white} or \ui{Viewport color}.
  \item[\ui{Renderer binary}] Path to the executable; the bare names
    \texttt{povray} and \texttt{tachyon} are resolved on \opt{PATH}. The
    path is remembered per engine.
\end{description}

Two ways to use it:

\begin{description}
  \item[\ui{Save Scene File\ldots}] Writes just the scene
    (\file{.pov} or \file{.dat}) for you to render elsewhere, or to edit by
    hand before rendering.
  \item[\ui{Render\ldots}] Writes the scene next to your chosen PNG output
    and invokes the renderer, streaming its output into the log pane so you
    can watch progress and diagnose failures.
\end{description}

\begin{calwarn}
Needs POV-Ray or Tachyon installed separately. If the binary cannot be
started, the log says so explicitly rather than failing silently.
\end{calwarn}

\section{Data exports at a glance}

Nearly every analysis tool exports its numbers, so figures can be
regenerated in your own plotting stack.

\begin{center}
\small
\begin{tabular}{@{}lll@{}}
\toprule
Source & Formats & Typical file \\
\midrule
RDF                     & CSV, DAT & \file{rdf.csv} \\
Bond length / angle     & CSV, DAT & \file{bond\_lengths.csv} \\
Structure factor        & CSV, DAT & \file{structure\_factor.csv} \\
XRD pattern             & CSV, DAT & \file{xrd\_pattern.csv} \\
XRD peak list           & CSV, DAT & \file{xrd\_peaks.csv} \\
Warren-Cowley           & CSV      & \file{warren\_cowley.csv} \\
Job metric series       & CSV, DAT & \file{energy.csv} \ldots \\
Band structure          & CSV, DAT & \file{bands.csv} \\
PDOS                    & CSV, DAT & \file{pdos.csv} \\
Isosurface mesh         & OBJ      & \file{isosurface.obj} \\
Volumetric slice        & CSV      & \file{slice.csv} \\
Phonon bands / DOS      & CSV      & \file{phonon\_bands.csv} \\
k-path                  & 5 codes  & Section~\ref{sec:bz} \\
ML datasets             & extxyz, ASE db & Section~\ref{sec:dataset-manager} \\
Brillouin zone figure   & PNG, SVG & \file{brillouin\_zone.svg} \\
\bottomrule
\end{tabular}
\end{center}
